\chapter{Arbitration}
\label{ch:arbitrering}

{\large\itshape How a bus without a master decides who gets to send.\par}

\begin{chapterbox}{In this chapter}
\begin{itemize}
  \item Multi-master: no node allocates the bus, and no node needs to.
  \item Arbitration bit by bit, worked by hand with two and three nodes.
  \item Why it is called non-destructive, and what that saves.
  \item Priority inversion, starvation, and how a real system avoids them.
\end{itemize}
\end{chapterbox}

\section{No master}

Most buses solve access with a master. One node decides who gets to talk, and the others wait
their turn. It works, and it costs a node that everyone depends on.

CAN has no such thing. Every node may start sending as soon as the bus has been idle for three bits
(the intermission). If two nodes start at the same time, the bus resolves it by itself, while it
happens, without anyone losing time.

The mechanism is the asymmetry \kapref{ch:buss} set out, used in a single place: during the
arbitration field, every transmitting node compares, bit by bit, what it sends with what is
actually on the bus.

\begin{aside}
\textbf{The rule, in full:} if a node sends \rec{} and reads back \dom{}, it has lost. It stops
sending immediately and becomes a receiver.
\end{aside}

If it sends dominant and reads dominant, all is as it should be --- either it is alone, or someone
else is sending the same bit. If it sends recessive and reads recessive, likewise.

The only combination that tells it anything is therefore recessive out, dominant in: someone else
is sending, and the other one has a zero where we have a one. Since the identifier is sent MSB
first, that means the other identifier is numerically \emph{lower}, and so has higher priority.

\section{Two nodes, bit by bit}

Node A wants to send the identifier \code{0x2A7}, node B \code{0x2A3}. Both start at the same time.

In binary, eleven bits, MSB first:

\begin{console}
bit:        1 2 3 4 5 6 7 8 9 10 11
A (0x2A7):  0 1 0 1 0 1 0 0 1  1  1
B (0x2A3):  0 1 0 1 0 1 0 0 0  1  1
                            ^ this is where they differ
\end{console}

\begin{booktable}{@{}lllll@{}}
\thead{Bit} & \thead{A sends} & \thead{B sends} & \thead{The bus} & \thead{Event} \\ \midrule
1--8 & same & same & same & Nobody notices anything \\
9 & \rec{} (1) & \dom{} (0) & \dom{} & A reads back dominant: \textbf{A loses} \\
10--11 & --- & sends on & B's bits & B continues undisturbed \\
\end{booktable}

Note what does \emph{not} happen: no collision, no lost time, no restart. B's frame continues from
bit 9 as if A had never existed. A's frame was never sent --- it just stopped, midway, and A waits
out B's frame to try again afterwards.

That is why it is called \textbf{non-destructive arbitration}: the conflict is resolved without
anyone's data being destroyed.

\bookfigure{arbitration}{Arbitration between two nodes. A withdraws and becomes a receiver at
bit 9; B's frame runs on uninterrupted.}{fig:arbitration}

\section{Three nodes}

With more nodes it is the same rule, applied by all of them at once. Three nodes with the
identifiers \code{0x100}, \code{0x0FF} and \code{0x2A7}:

\begin{console}
bit:        1 2 3 4 5 6 7 8 9 10 11
A (0x100):  0 0 1 0 0 0 0 0 0  0  0
B (0x0FF):  0 0 0 1 1 1 1 1 1  1  1
C (0x2A7):  0 1 0 1 0 1 0 0 1  1  1
              ^ ^
              | bit 3: A sends recessive, B dominant -> A drops out
              bit 2: C sends recessive, A and B dominant -> C drops out
\end{console}

At bit 1 all three send dominant, so nobody notices anything. At bit 2 C drops out, at bit 3 A
drops out. B wins, and has the lowest identifier --- which is always the outcome, because ``lowest
numerically'' and ``first dominant bit'' are the same thing when the bits are sent MSB first.

\section{What the loser does}

The losing node becomes a receiver for the rest of the frame. That is not just politeness: it
\emph{must} receive the winning frame, since it may very well be meant for it.

When the bus has become idle again, it tries again. And since its identifier is unchanged, it will
lose again if the same high-priority node sends again --- which leads to the next section.

\begin{simplification}
\note[Simplification] The simplified controller in \kapref{ch:kursen} does \emph{not} receive the
frame it lost to. It leaves the transmitter role, sets its error flag and goes back to idle. That
is a deliberate simplification, and it is visible to the software: a node that has lost an
arbitration misses the frame that won.
\end{simplification}

\section{Starvation and priority inversion}

Arbitration is fair to priority, not to nodes. A node with a high identifier on a heavily loaded
bus can in principle never get through --- it loses every time.

This is known and accepted behaviour, and it is handled in the system design rather than in the
protocol:

\begin{itemize}
  \item \textbf{Bus load is kept down.} The rule of thumb in vehicle systems is not to plan for
    more than about 30--50~\% load, precisely to leave room for low-priority frames.
  \item \textbf{Identifiers are planned.} The priority order is set by what is time-critical, not
    by what happened to get which number.
  \item \textbf{Transmissions are made periodic.} A node that sends as often as it can is a node
    that starves everyone below it.
\end{itemize}

\begin{aside}
\note A frame that has lost arbitration is delayed, but not corrupted. That is why CAN is called
\emph{deterministic} for the highest-priority frame: its worst-case delay can be calculated. For the
other frames it can too, but the answer grows the further down the order they are.
\end{aside}

\section{What the transmitter never learns}

A node that has lost arbitration knows that it lost. A node that won only knows that it got
through.

Above all: \textbf{the transmitter never learns whether anyone lost to it.} There is no counter, no
report, no difference in the frame. That does not matter to CAN, but it matters to whoever is
debugging a system where one node never seems to get through --- the symptom is visible only at the
loser.

\section*{Exercises}

\exercise{Two nodes}{Calculation}
\label{ex:4:1}
Node A sends the identifier \code{0x155}, node B sends \code{0x154}. Both start at the same time.
\begin{enumerate}[a)]
  \item Write out both identifiers in binary, eleven bits, MSB first.
  \item At which bit do they differ?
  \item Who wins, and what does the loser do?
\end{enumerate}

\exercise{Three nodes}{Calculation}
\label{ex:4:2}
Three nodes start sending at the same time, with the identifiers \code{0x400}, \code{0x200} and
\code{0x201}.
\begin{enumerate}[a)]
  \item At which bit does the first node drop out?
  \item At which bit is it settled between the remaining two?
  \item Who wins?
\end{enumerate}

\exercise{The same identifier}{Reasoning}
\label{ex:4:3}
Two nodes by mistake send the same identifier at the same time, but with different data.
\begin{enumerate}[a)]
  \item What happens during the arbitration field?
  \item What happens when the data fields start to differ?
  \item Why is this a design error rather than an operational fault?
\end{enumerate}

\exercise{Starvation}{Reasoning}
\label{ex:4:4}
A node with the identifier \code{0x7FF} never gets through on a heavily loaded bus.
\begin{enumerate}[a)]
  \item Why not?
  \item Is it a flaw in the protocol? Justify your answer.
  \item Give two measures at system level that would help.
\end{enumerate}
